% Reversi_All.tex - Combined source listing for Lab_1_Reversi_Board/myCode
\documentclass[11pt]{article}
\usepackage[utf8]{inputenc}
\usepackage[a4paper,margin=1in]{geometry}
\usepackage{xcolor}
\usepackage{courier}
\usepackage{listings}
\usepackage{hyperref}

% Color scheme
\definecolor{kwcolor}{RGB}{0,0,180}
\definecolor{commentcolor}{RGB}{0,120,0}
\definecolor{stringcolor}{RGB}{180,0,0}
\definecolor{bg}{RGB}{248,248,248}

\lstset{
  backgroundcolor=\color{bg},
  language=C++,
  basicstyle=\footnotesize\ttfamily,
  keywordstyle=\color{kwcolor}\bfseries,
  commentstyle=\color{commentcolor}\itshape,
  stringstyle=\color{stringcolor},
  numberstyle=\tiny\color{gray},
  numbers=left,
  stepnumber=1,
  numbersep=8pt,
  tabsize=4,
  breaklines=true,
  showstringspaces=false,
  frame=single,
  captionpos=b
}

\begin{document}

\begin{center}
  {\LARGE Reversi - Combined Source Listing}
  \\[6pt]
    % file list removed per request
\end{center}
\vspace{1pt}

\section*{ReversiBoard.h}
\begin{lstlisting}
/* ReversiBoard.h - header (short notes) */

#ifndef REVERSIBOARD_H_
#define REVERSIBOARD_H_

// FieldState: enum to represent a cell's contents (BLACK, WHITE, BLANK)
// Exam note: enums are integer-like named constants; easier to read than raw ints.
enum FieldState
{
	BLACK, ///< Represents a black piece.
	WHITE, ///< Represents a white piece.
	BLANK  ///< Represents an empty field.
};

// ReversiBoard: holds 8x8 board. OOP: groups data + methods.
class ReversiBoard
{
private:
	// No private members

public:
	// 2D array representing the board with 8x8 fields.
	FieldState board[8][8];

	// Constructor: sets up initial board pieces.
	// Constructor/Destructor note: constructors initialize objects; destructors
	// free resources (none needed here).
	ReversiBoard();

	// Destructor (no special cleanup required for this simple class).
	~ReversiBoard();

	// Returns the state of the cell at (column, row).
	// Important: parameters are (column, row) — check order when calling.
	FieldState getFieldState(int column, int row);

	// Attempts to place `piece` at (column, row). Returns true if legal move
	// and flips opponent pieces accordingly. Uses loops and conditionals.
	// Exam note: this function demonstrates bounds checks, loops, and
	// modifying a 2D array (board).
	bool setFieldState(int column, int row, FieldState piece);

};

#endif /* REVERSIBOARD_H_ */
\end{lstlisting}

\section*{ReversiBoard.cpp}
\begin{lstlisting}
/* ReversiBoard.cpp
 * Simple implementation notes and exam hints added near important lines.
 * Learning: shows loops, conditionals, array access, and basic OOP usage.
 */

#include <iostream>
#include "ReversiBoard.h"

// Constructor: initialize board cells and place starting pieces.
// Syntax note: nested loops used to initialize 2D array.
ReversiBoard::ReversiBoard()
{
    // Initialize all fields to BLANK (looping over columns and rows).
    for (int column = 0; column < 8; column++)
    {
        for (int row = 0; row < 8; row++)
        {
            board[column][row] = FieldState::BLANK; // array access: board[c][r]
        }
    }

    // Starting position: center pieces (indices 3 and 4)
    board[3][3] = BLACK;
    board[4][4] = BLACK;
    board[3][4] = WHITE;
    board[4][3] = WHITE;
}

// Destructor: no dynamic memory, so left empty.
ReversiBoard::~ReversiBoard() {}

// Returns state at (column,row)
FieldState ReversiBoard::getFieldState(int column, int row)
{
    return board[column][row];
}

// Place piece and flip opponent pieces; returns true if legal
bool ReversiBoard::setFieldState(int column, int row, FieldState piece)
{
    // If cell is not empty, move is invalid.
    if (getFieldState(column, row) != BLANK)
    {
        return false; // condition check example
    }

    bool isAllowedToPlace = false; // flag set to true when a valid flip is found

    // EAST: move right from placed piece
    for (int col = column + 1; col < 8; col++)
    {
        if (getFieldState(col, row) == BLANK) // empty -> no flanking
            break;

        if (getFieldState(col, row) == piece)
        {
            if (col == column + 1) // adjacent same-color piece -> nothing to flip
                break;

            // Flip pieces between `column` and `col` (inclusive of boundaries)
            for (int j = column; j < col; j++)
            {
                board[j][row] = piece; // array write
            }
            isAllowedToPlace = true;
            break;
        }
    }

    // WEST
    for (int col = column - 1; col >= 0; col--)
    {
        if (getFieldState(col, row) == BLANK)
            break;

        if (getFieldState(col, row) == piece)
        {
            if (col == column - 1)
                break;

            for (int j = column; j > col; j--)
            {
                board[j][row] = piece;
            }
            isAllowedToPlace = true;
            break;
        }
    }

    // SOUTH
    for (int row1 = row + 1; row1 < 8; row1++)
    {
        if (getFieldState(column, row1) == BLANK)
            break;

        if (getFieldState(column, row1) == piece)
        {
            if (row1 == row + 1)
                break;

            for (int j = row; j < row1; j++)
            {
                board[column][j] = piece;
            }
            isAllowedToPlace = true;
            break;
        }
    }

    // NORTH
    for (int row1 = row - 1; row1 >= 0; row1--)
    {
        if (getFieldState(column, row1) == BLANK)
            break;

        if (getFieldState(column, row1) == piece)
        {
            if (row1 == row - 1)
                break;

            for (int j = row; j > row1; j--)
            {
                board[column][j] = piece;
            }
            isAllowedToPlace = true;
            break;
        }
    }

    // NORTH-EAST diagonal
    for (int i = 1; i < 8; i++)
    {
        if (!(column >= 0 && column < 8 && row >= 0 && row < 8))
            break; // safety check for indices

        if (getFieldState(column + i, row - i) == BLANK)
            break;

        if (getFieldState(column + i, row - i) == piece)
        {
            if (i == 1)
                break;

            for (int j = 0; j < i; j++)
            {
                board[column + j][row - j] = piece;
            }
            isAllowedToPlace = true;
            break;
        }
    }

    // NORTH-WEST diagonal
    for (int i = 1; i < 8; i++)
    {
        if (!(column >= 0 && column < 8 && row >= 0 && row < 8))
            break;

        if (getFieldState(column - i, row - i) == BLANK)
            break;

        if (getFieldState(column - i, row - i) == piece)
        {
            if (i == 1)
                break;

            for (int j = 0; j < i; j++)
            {
                board[column - j][row - j] = piece;
            }
            isAllowedToPlace = true;
            break;
        }
    }

    // SOUTH-EAST diagonal
    for (int i = 1; i < 8; i++)
    {
        if (!(column >= 0 && column < 8 && row >= 0 && row < 8))
            break;

        if (getFieldState(column + i, row + i) == BLANK)
            break;

        if (getFieldState(column + i, row + i) == piece)
        {
            if (i == 1)
                break;

            for (int j = 0; j < i; j++)
            {
                board[column + j][row + j] = piece;
            }
            isAllowedToPlace = true;
            break;
        }
    }

    // SOUTH-WEST diagonal
    for (int i = 1; i < 8; i++)
    {
        if (!(column >= 0 && column < 8 && row >= 0 && row < 8))
            break;

        if (getFieldState(column - i, row + i) == BLANK)
            break;

        if (getFieldState(column - i, row + i) == piece)
        {
            if (i == 1)
                break;

            for (int j = 0; j < i; j++)
            {
                board[column - j][row + j] = piece;
            }
            isAllowedToPlace = true;
            break;
        }
    }

    // If any direction produced flips, place the new piece.
    if (isAllowedToPlace)
    {
        board[column][row] = piece;
    }

    return isAllowedToPlace; // function returns boolean result
}
\end{lstlisting}

\section*{ReversiConsoleView.h}
\begin{lstlisting}
/* ReversiConsoleView.h
 * Simple notes: console view for printing the board.
 * Learning: demonstrates pointer member to another object (ReversiBoard*).
 */

#ifndef REVERSICONSOLEVIEW_H_
#define REVERSICONSOLEVIEW_H_

#include "ReversiBoard.h"

// ReversiConsoleView: prints the board to the console.
// Pointer note: stores `ReversiBoard *reversiBoard` — shows simple pointer use.
class ReversiConsoleView
{
private:
	/// Pointer to the ReversiBoard object to be displayed.
	ReversiBoard *reversiBoard;

public:
	// Constructor: takes pointer to a ReversiBoard to display.
	// Exam note: passing pointer avoids copying the whole board; caller owns board.
	ReversiConsoleView(ReversiBoard *reversiBoard);

	// Destructor: no ownership cleanup here (does not delete the board pointer).
	~ReversiConsoleView();

	// Prints board to console using simple characters (O/X/+).
	// Loops and conditionals demo: iterates rows/columns and branches on state.
	void print();
};

#endif /* REVERSICONSOLEVIEW_H_ */
\end{lstlisting}

\section*{ReversiConsoleView.cpp}
\begin{lstlisting}
/* ReversiConsoleView.cpp
 * Simple implementation notes and exam hints added for printing logic.
 * Shows pointer member access and basic I/O with std::cout/std::cin.
 */

#include <iostream>
#include "ReversiConsoleView.h"
#include "ReversiBoard.h"

using namespace std;

// Constructor: store pointer to board
ReversiConsoleView::ReversiConsoleView(ReversiBoard *reversiBoard)
{
	this->reversiBoard = reversiBoard; // pointer assignment
}

// Destructor: nothing to free
ReversiConsoleView::~ReversiConsoleView() {}

// Print board: nested loops, map enum->symbol
void ReversiConsoleView::print()
{
	cout << " 0  1  2  3  4  5  6  7 " << endl; // column headers

	for (int row = 0; row < 8; row++)
	{
		for (int column = 0; column < 8; column++)
		{
			// Map FieldState to symbol: exam-friendly shorthand
			if (reversiBoard->board[column][row] == WHITE)
			{
				cout << " O "; // WHITE -> O (circle)
			}
			else if (reversiBoard->board[column][row] == BLACK)
			{
				cout << " X "; // BLACK -> X (cross)
			}
			else
			{
				cout << " + "; // BLANK -> + (empty)
			}

			if (column == 7)
				cout << endl << endl; // new line after each row
		}
	}
}
\end{lstlisting}

\section*{main.cpp}
\begin{lstlisting}
/* main.cpp - simple runner */

#include <iostream>
#include <cstdlib>
#include "ReversiBoard.h"
#include "ReversiConsoleView.h"

using namespace std;

int main()
{
	// Disable output buffering so console updates appear immediately.
	setvbuf(stdout, NULL, _IONBF, 0);

    // Create board and console view (pass pointer)
	ReversiBoard myBoard;
	ReversiConsoleView myConsoleView(&myBoard); // pass pointer to avoid copying

	int column, row;
    bool isWhitesTurn = true; // whose turn

    // Main game loop
	do
	{
		if (isWhitesTurn)
		{
			myConsoleView.print();
			cout << "White's Turn" << endl;
			cout << "Enter Column: ";
            cin >> column; // input
			cout << "Enter Row: ";
			cin >> row;

			// If move valid, swap turn (example of function call and boolean result)
			if (myBoard.setFieldState(column, row, WHITE))
			{
				isWhitesTurn = false; // switch player
			}
		}
		else
		{
			myConsoleView.print();
			cout << "Black's Turn" << endl;
			cout << "Enter Column: ";
			cin >> column;
			cout << "Enter Row: ";
			cin >> row;

			if (myBoard.setFieldState(column, row, BLACK))
			{
				isWhitesTurn = true;
			}
		}
    } while (true); // Infinite loop (no end condition here)

	return 0;
}
\end{lstlisting}

\end{document}
